\section{One tail heavier: the dominated regime}\label{sec:dominated}

Throughout this section (DOM) holds. The plus runs have the heavier tail,
$T_\minus=o(T_\plus)$, with $\alpha_\plus\in(0,1)$, and the minus runs have an index
$\alpha_\minus\in(0,1)\cup(1,2)$, so their mean $\mu_\minus=\E X_\minus$ is finite exactly when
$\alpha_\minus>1$. We assume (RV) and (SPAN) for both signs. Then $A_N/N\to1$ in probability,
and the object is the time $N-A_N$ spent on the lighter side. The plus side alone sets the
number of cycles. By time $N$ about $E_\plus/T_\plus(N)$ plus runs are complete, with
$E_\plus$ the Mittag-Leffler variable of Section~\ref{sec:elementary}, and the minus time is
the sum of that many minus runs. So the scale and the limit are
\begin{align*}
 \mu_\minus<\infty:&\qquad D_N=\frac{\mu_\minus}{T_\plus(N)},\qquad \rho=\phi_\plus,\qquad \omega(j)=1,\\
 \alpha_\minus<1:&\qquad D_N=a_\minus\bigl(1/T_\plus(N)\bigr),\qquad \rho=\text{the density of }E_\plus^{1/\alpha_\minus}Z_\minus,\qquad \omega=\omega_\minus,
\end{align*}
where $Z_\minus$ is the positive stable variable with L\'evy tail $z^{-\alpha_\minus}$,
independent of $E_\plus$, so that $T_\minus(D_N)\sim T_\plus(N)$ in the second case. In both
cases $D_N$ is nondecreasing and regularly varying in $N$, with index $\alpha_\plus$ when
$\mu_\minus<\infty$ and $\alpha_\plus/\alpha_\minus$ when $\alpha_\minus<1$, and $D_N=o(N)$, because
$NT_\plus(N)\to\infty$ in the first case and $T_\minus(D_N)\sim T_\plus(N)\gg T_\minus(N)$ in the
second, so $T_\plus(N-j)\sim T_\plus(N)$ uniformly in $j\le KD_N$. We write $\phi=\phi_\plus$,
and we put
\[
 \Psi(y)=\rho(y)\quad(\mu_\minus<\infty),\qquad
 \Psi(y)=\rho(y)\Big/\Bigl(\frac{\sin\pi\alpha_\minus}\pi\,y^{\alpha_\minus-1}\Bigr)\quad(\alpha_\minus<1),
\]
so that $\Psi(0+)=1$ in both cases. For $\mu_\minus<\infty$ this is $\phi(0)=1$. For
$\alpha_\minus<1$, the substitution $x=(y/v)^{\alpha_\minus}$ in
$\rho(y)=\int_0^\infty\phi(x)x^{-1/\alpha_\minus}g_\minus(yx^{-1/\alpha_\minus})\,dx$ gives
$\rho(y)=\alpha_\minus y^{\alpha_\minus-1}\int_0^\infty\phi((y/v)^{\alpha_\minus})g_\minus(v)v^{-\alpha_\minus}dv$,
and as $y\to0$ dominated convergence gives
$\rho(y)\sim\alpha_\minus y^{\alpha_\minus-1}\E Z_\minus^{-\alpha_\minus}=\frac{\sin\pi\alpha_\minus}\pi y^{\alpha_\minus-1}$.
When $\alpha_\plus=\alpha_\minus=\alpha$ and $b=0$ the limit is the ratio $Z_\minus/Z_\plus$ of two
independent copies of the same stable law, with the explicit density
\begin{equation}\label{eq:ratio-density}
 \rho(y)=\frac{\sin\pi\alpha}\pi\,\frac{y^{\alpha-1}}{1+2y^\alpha\cos\pi\alpha+y^{2\alpha}},
\end{equation}
which is also the limit of $b^{1/\alpha}\La(1-b^{1/\alpha}y)$ as $b\to0$. So the balanced and
the dominated regime join continuously at $b=0$. We also use the condition
\begin{equation}\label{eq:W}\tag{W$_\plus$}
 D_nT_\plus(n)\sum_{m\le\delta/T_\plus(n)}u^\plus_m(n)\to0\qquad(n\to\infty)\quad\text{for some }\delta>0 .
\end{equation}
For $\delta<\delta'$ the extra terms number at most $\delta'/T_\plus(n)$ and each is at most
$C/a_\plus(\delta/T_\plus(n))\le C_\delta/n$ by Lemma~\ref{lem:cd}(i), so the difference is
$O(D_n/n)\to0$, and (W$_\plus$) for one $\delta$ is (W$_\plus$) for every $\delta$.

\begin{introtheorem}[the dominated regime]\label{thm:dominated}
Assume (DOM) and (SPAN), and let the start be fresh with either sign.
\begin{enumerate}[(a)]
\item (The global law.) $(N-A_N)/D_N$ converges in law to the law with density $\rho$.
\item (The lower half.) For $0<y_1<y_2<\infty$,
$\liminf_N\inf_{y_1D_N\le j\le y_2D_N}\bigl(D_NP_N(N-j)-\rho(j/D_N)\bigr)\ge0$.
\item (The local law.) If $\mu_\minus<\infty$ or $\alpha_\minus>1/2$, then
\[
 \sup_{y_1D_N\le j\le y_2D_N}\bigl|D_NP_N(N-j)-\rho(j/D_N)\bigr|\to0 .
\]
If $\alpha_\minus\le1/2$, this holds if and only if (SRT$_\minus$) and (W$_\plus$) hold.
(W$_\plus$) holds under (SRT$_\plus$) and whenever $\alpha_\minus>1/2$.
\item (The end layers.) For each fixed $j\ge0$, $P_N(N-j)/T_\plus(N)\to\kappa^{(\sgn_0)}_\plus(j)$
of~\eqref{eq:kappa-fresh}, under (RV) for the plus sign alone, and for each fixed $i\ge0$,
$P_N(i)/T_\minus(N)\to\kappa^{(\sgn_0)}_\minus(i)$ under (RV) for the minus sign alone. Under the
condition of (c), and if moreover $\alpha_\minus\ne1/2$, uniformly over $0\le j\le y_2D_N$,
\[
 P_N(N-j)=T_\plus(N)\bigl[\kappa^{(\sgn_0)}_\plus(j)\,\Psi(j/D_N)+o(\omega(j))\bigr].
\]
At $\alpha_\minus=1/2$ the same holds under (SRT$_\plus$) with $\alpha_\plus<1/2$.
\end{enumerate}
\end{introtheorem}

Part (a) is known in distribution. For $\mu_\minus<\infty$ it is a Darling--Kac
law~\cite{darlingkac1957} for the time spent in a part of a null recurrent system that is
left at a positive rate, and for the lattice walk it is stated without proof by C\'enac, Le
Ny, de Loynes and Offret~\cite[Section~4.3]{clnlo2019}. Yamato~\cite{yamato2026} proves
limits of this type for processes time-changed by inverse subordinators. Lamperti's
theorem~\cite{lamperti1958} covers unequal tails but gives only $A_N/N\to1$ there. We include
(a) because it follows from (b) in two lines. Parts (b) to (d) are local, and we did not find
them in the literature. The local law needs two inputs only: the local limit theorem for the
lighter sums and the local form of the inverse renewal count of the heavier runs, which is
Lemma~\ref{lem:Q}. The strong renewal theorem enters through the lighter law alone, and only
when its mean is infinite with index at most $1/2$. In that case the heavier law must also
satisfy (W$_\plus$). We do not know whether (W$_\plus$) can fail when $\alpha_\plus<1/2$ and
(SRT$_\plus$) fails.

Throughout the proofs we write $n=N-j$ and, by Proposition~\ref{prop:fresh-identity} for a
fresh plus start,
\[
 P_N(N-j)=\mathrm I_N(j)+\mathrm R_N(j),\qquad
 \mathrm I_N(j)=\sum_{m\ge1}u^\minus_m(j)\,q^\plus_{m+1}(n),\qquad
 \mathrm R_N(j)=\sum_{m\ge1}u^\plus_m(n)\,q^\minus_m(j).
\]
For a fresh minus start the kernels are $q^\plus_m$ and $q^\minus_{m+1}$, which changes no
bound and no limit below. We first record where the mass of $u_\minus(j)$ sits in $m$.

\begin{lemma}[the index of the minus runs]\label{lem:R}
Assume (RV) and (SPAN) for the minus sign.
\begin{enumerate}[(i)]
\item If $\mu_\minus<\infty$, then
$\lim_{A\to\infty}\limsup_{j\to\infty}\sum_{m:\ |m-j/\mu_\minus|>Aa_\minus(j)}u^\minus_m(j)=0$, where
$a_\minus$ is the scale of the minus sums.
\item If $\alpha_\minus<1$, then for every $\epsilon>0$ there is $C_\epsilon$ with
$\sum_{m>K/T_\minus(j)}u^\minus_m(j)\le C_\epsilon K^{1-1/\alpha_\minus+\epsilon}/(jT_\minus(j))$ for all
$K\ge1$ and $j\ge1$.
\end{enumerate}
\end{lemma}

\begin{proof}
(i) By the local limit theorem for the minus sums~\cite[Lemma~A.1]{paperF},
$u^\minus_m(j)=a_\minus(m)^{-1}[g_\minus((j-m\mu_\minus)/a_\minus(m))+o(1)]$ uniformly in $j$, so the
sum over $|m-j/\mu_\minus|\le Aa_\minus(j)$ is a Riemann sum that converges to
$\mu_\minus^{-1}\int_{|z|\le A'}g_\minus(z)\,dz$ with $A'\to\infty$ as $A\to\infty$, while $u_\minus(j)\to1/\mu_\minus$ by the renewal
theorem. (ii) By Lemma~\ref{lem:cd}(i), $u^\minus_m(j)\le C/a_\minus(m)$, and since $1/a_\minus(m)$ is
regularly varying with index $-1/\alpha_\minus<-1$, $\sum_{m>M}1/a_\minus(m)\le CM/a_\minus(M)$. With
$M=K/T_\minus(j)$, Potter's bounds for the regularly varying $a_\minus$ give
$a_\minus(M)\ge c_\epsilon K^{1/\alpha_\minus-\epsilon}a_\minus(1/T_\minus(j))\ge c_\epsilon K^{1/\alpha_\minus-\epsilon}j$
uniformly in $K\ge1$, which gives the bound.
\end{proof}

\subsection{The bulk: proof of parts (a), (b) and (c)}\label{sec:dom-bulk}

Let $j=yD_N$ with $y\in[y_1,y_2]$, so that $T_\plus(n)\sim T_\plus(N)$ uniformly.

\emph{The main term, $\mu_\minus<\infty$.} On the window $|m-j/\mu_\minus|\le Aa_\minus(j)$ we have
$mT_\plus(n)=y+O(Aa_\minus(j)/D_N)+o(1)=y+o(1)$, since $a_\minus(j)/j\to0$ when $\alpha_\minus>1$. So
Lemma~\ref{lem:Q}(ii) and the uniform continuity of $\phi$ give
$q^\plus_{m+1}(n)=T_\plus(N)[\phi(y)+o(1)]$ on the window, and by Lemma~\ref{lem:R}(i) the weights
$u^\minus_m(j)$ on the window add up to $1/\mu_\minus+o_A(1)$. Off the window
Lemma~\ref{lem:key} gives $q^\plus_{m+1}(n)\le CT_\plus(N)$, and the weights add up to $o_A(1)$.
Hence $\mathrm I_N(j)=(T_\plus(N)/\mu_\minus)(\phi(y)+o(1))=(\phi(y)+o(1))/D_N$, with no condition
beyond (RV) and (SPAN).

\emph{The main term, $\alpha_\minus<1$.} Now $T_\minus(j)\sim y^{-\alpha_\minus}T_\plus(N)$. Split
the sum over $m$ at $\varepsilon/T_\plus(N)$ and $K/T_\plus(N)$. On the middle range,
Lemma~\ref{lem:cd}(i) gives $u^\minus_m(j)=a_\minus(m)^{-1}[g_\minus(j/a_\minus(m))+o(1)]$ and
Lemma~\ref{lem:Q}(ii) gives $q^\plus_{m+1}(n)=T_\plus(N)[\phi(mT_\plus(N))+o(1)]$, both uniformly.
Since $a_\minus(x/T_\plus(N))\sim x^{1/\alpha_\minus}D_N$ uniformly for $x\in[\varepsilon,K]$, the sum
is a Riemann sum with mesh $T_\plus(N)$ in $x=mT_\plus(N)$, and
\[
 T_\plus(N)\sum_m\frac1{a_\minus(m)}\,g_\minus\Bigl(\frac j{a_\minus(m)}\Bigr)\phi\bigl(mT_\plus(N)\bigr)
 \to\frac1{D_N}\int_\varepsilon^K\phi(x)\,x^{-1/\alpha_\minus}g_\minus\bigl(yx^{-1/\alpha_\minus}\bigr)\,dx ,
\]
uniformly in $y$, and the integral tends to $\rho(y)$ as $\varepsilon\to0$ and $K\to\infty$. On the
large range Lemma~\ref{lem:key} and Lemma~\ref{lem:R}(ii) give at most
$CT_\plus(N)K^{1-1/\alpha_\minus+\epsilon}/(jT_\minus(j))\le CK^{1-1/\alpha_\minus+\epsilon}/D_N$, with a fixed $\epsilon<1/\alpha_\minus-1$, since
$jT_\minus(j)\sim y^{1-\alpha_\minus}D_NT_\plus(N)$. On the small range Lemma~\ref{lem:key} gives at most
\begin{equation}\label{eq:dom-small}
 CT_\plus(N)\sum_{m<\varepsilon'/T_\minus(j)}u^\minus_m(j)=\frac{CT_\plus(N)}{jT_\minus(j)}\,I_\minus(\varepsilon',j),
\end{equation}
with $I_\minus$ as in~\eqref{eq:cd41} and $\varepsilon'=C\varepsilon$. The prefactor is of order $1/D_N$,
and by Lemma~\ref{lem:cd}(iv) the factor $I_\minus(\varepsilon',j)$ tends to $0$ as $j\to\infty$ and then
$\varepsilon'\to0$ if and only if (SRT$_\minus$) holds.

\emph{The second term, $\alpha_\minus<1$.} By Lemma~\ref{lem:key} for the minus sign,
$q^\minus_m(j)\le CT_\minus(j)\le CT_\plus(N)$. For $m>\delta/T_\plus(n)$, Lemma~\ref{lem:cd}(i) gives
$u^\plus_m(n)\le C/a_\plus(\delta/T_\plus(n))\le C_\delta/N$, and $\sum_mq^\minus_m(j)=1$, so this part
is $O(1/N)=o(1/D_N)$. The part $m\le\delta/T_\plus(n)$ is at most
$CT_\plus(N)\sum_{m\le\delta/T_\plus(n)}u^\plus_m(n)$, which is $o(1/D_N)$ exactly under (W$_\plus$),
since $D_n\sim D_N$. Under (SRT$_\plus$), $\sum_mu^\plus_m(n)=u_\plus(n)\le C/(nT_\plus(n))$, so the
left side of (W$_\plus$) is at most $CD_n/n\to0$. In general Lemma~\ref{lem:cd}(ii) gives
$\sum_{m\le M}u^\plus_m(n)\le CT_\plus(n)\sum_{m\le M}m/a_\plus(m)$ with $M=\delta/T_\plus(n)$. The
sequence $m/a_\plus(m)$ is regularly varying with index $1-1/\alpha_\plus$. If $\alpha_\plus<1/2$ the
series converges and the sum is at most $CT_\plus(n)$, so the left side of (W$_\plus$) is at most
$CD_nT_\plus(n)^2\sim CD_nT_\minus(D_n)^2$, which tends to $0$ when $\alpha_\minus>1/2$, because
$xT_\minus(x)^2$ is regularly varying with index $1-2\alpha_\minus<0$. If $\alpha_\plus\ge1/2$ the partial
sums are at most $CM^2\lambda(M)/a_\plus(M)$ with $\lambda$ slowly varying ($\lambda=1$ for
$\alpha_\plus>1/2$, by Karamata's theorem~\cite[Proposition~1.5.8 and~1.5.9a]{bgt1989}), and
$a_\plus(\delta/T_\plus(n))\ge c_\delta n$, so the sum is at most $C_\delta\lambda(1/T_\plus(n))/(nT_\plus(n))$
and the left side of (W$_\plus$) is at most $C_\delta\lambda(1/T_\plus(n))D_n/n$. Under (DOM),
$\alpha_\plus\ge1/2$ forces $\alpha_\minus\ge\alpha_\plus\ge1/2$, and $\alpha_\minus>1/2$ leaves two cases.
If $\alpha_\plus<\alpha_\minus$, then $D_n/n$ is regularly varying with a negative index, and the
bound tends to $0$. If $\alpha_\plus=\alpha_\minus>1/2$, then $\lambda\equiv1$ and $D_n/n\to0$. This proves
the last sentence of (c).

\emph{The second term, $\mu_\minus<\infty$.} Here $q^\minus_m(j)$ is not $O(T_\minus(j))$ uniformly in
$m$, and we count directly. Split at $m=\delta j$ with $\delta\mu_\minus\le1/4$. For $m>\delta j$,
Lemma~\ref{lem:cd}(i) gives $u^\plus_m(n)\le C/a_\plus(\delta j)$, and $a_\plus(\delta j)\ge cN$ because
$j$ is of order $1/T_\plus(N)$, so with $\sum_mq^\minus_m(j)=1$ this part is $O(1/N)=o(T_\plus(N))$.
For $m\le\delta j$ write $q^\minus_m(j)=\sum_{r<j}u^\minus_{m-1}(r)T_\minus(j-r)$. The part $r<j/2$ is
at most $T_\minus(j/2)\le CT_\minus(j)$. For $r\ge j/2$ we have $r-(m-1)\mu_\minus\ge j/4$, so the
local large deviation bound for the minus sums~\cite[Lemma~A.1(iii)]{paperF} gives
$u^\minus_{m-1}(r)\le C(m/a_\minus(m))T_\minus(j)$, and $\sum_rT_\minus(j-r)\le\mu_\minus$. Hence
$q^\minus_m(j)\le CT_\minus(j)(1+m/a_\minus(m))$ for $m\le\delta j$. With
$u^\plus_m(n)\le C(m/a_\plus(m))T_\plus(n)$ from Lemma~\ref{lem:cd}(ii), the part $m\le\delta j$ of
$\mathrm R_N(j)$ is at most
\begin{multline*}
 CT_\plus(N)\,T_\minus(j)\sum_{m\le\delta j}\frac m{a_\plus(m)}\Bigl(1+\frac m{a_\minus(m)}\Bigr)\\
 \le CT_\plus(N)\,T_\minus(j)\,\lambda(j)\Bigl[\max\Bigl(1,\frac{j^2}{a_\plus(j)}\Bigr)+\max\Bigl(1,\frac{j^3}{a_\plus(j)a_\minus(j)}\Bigr)\Bigr]
\end{multline*}
with $\lambda$ slowly varying. The two products with $T_\minus(j)$ are regularly varying in $j$
with indices $-\alpha_\minus+\max(0,2-1/\alpha_\plus)$ and
$-\alpha_\minus+\max(0,3-1/\alpha_\plus-1/\alpha_\minus)$, and both are negative, since
$2-1/\alpha_\plus<1<\alpha_\minus$ and $3-1/\alpha_\plus-1/\alpha_\minus-\alpha_\minus<2-\alpha_\minus-1/\alpha_\minus\le0$.
So $\mathrm R_N(j)=o(T_\plus(N))=o(1/D_N)$ as $j\to\infty$, with no condition on $j/D_N$. In the
mixed region the bulk sits at the height $T_\plus(N)$ of the end layer, and a minus run in
progress at time $N$ costs a polynomial factor against it.

\emph{Part (b).} Every term is nonnegative. Dropping the small range of $\mathrm I_N$ and all of
$\mathrm R_N$ leaves the middle and large ranges, which give
$D_NP_N(N-j)\ge\rho(y)-\eta(\varepsilon,K)+o(1)$ uniformly, with $\eta(\varepsilon,K)\to0$ as
$\varepsilon\to0$ and $K\to\infty$.

\emph{Part (c), sufficiency.} If $\mu_\minus<\infty$ the two terms give the claim with no
condition. If $\alpha_\minus<1$, under (SRT$_\minus$) and (W$_\plus$) the small range
of~\eqref{eq:dom-small} and the second term are $o(1/D_N)$, and the rest converges. If
$\alpha_\minus>1/2$ both conditions are automatic, by Lemma~\ref{lem:cd}(iii) and the last
paragraph on (W$_\plus$).

\emph{Part (c), necessity for $\alpha_\minus\le1/2$.} Suppose (SRT$_\minus$) fails. Since
$I_\minus(\varepsilon',j)$ is nondecreasing in $\varepsilon'$, Lemma~\ref{lem:cd}(iv) gives $c>0$ such
that for every $\varepsilon'>0$, $I_\minus(\varepsilon',j)\ge c$ along a sequence of $j\to\infty$ that may
depend on $\varepsilon'$. Fix $y\in(y_1,y_2)$ and a small $\eta>0$. For each such $j$ choose $N$ with
$j/D_N\in[y-\eta,y+\eta]$, which is possible because $D_N$ is nondecreasing and
$D_{N+1}/D_N\to1$. By Lemma~\ref{lem:Q}(i), $q^\plus_{m+1}(n)\ge\frac12T_\plus(N)$ for
$mT_\plus(n)\le\delta_0$, which covers $m<\varepsilon'/T_\minus(j)$ once $\varepsilon'$ is small. So the
small range of $\mathrm I_N(j)$ is at least $\frac12T_\plus(N)\sum_{m<\varepsilon'/T_\minus(j)}u^\minus_m(j)\ge c'/D_N$,
with $c'>0$ independent of $\varepsilon'$, while the middle range is at least
$(\rho(y)-\eta(\varepsilon',K)-o(1))/D_N$. Choosing $\varepsilon'$ and $K$ so that
$\eta(\varepsilon',K)<c'/4$, and $\eta$ so that $|\rho(j/D_N)-\rho(y)|<c'/4$, gives
$D_NP_N(N-j)-\rho(j/D_N)\ge c'/4$ along the sequence, so the local law fails. Suppose instead
that (W$_\plus$) fails, so that $D_nT_\plus(n)\sum_{m\le\delta/T_\plus(n)}u^\plus_m(n)\ge c$ along a
sequence of $n$, for a $\delta$ so small that $mT_\minus(j)\le\delta_0$ on that range for
$j\ge y_1D_N$. Lemma~\ref{lem:Q}(i) for the minus sign gives $q^\minus_m(j)\ge\frac12T_\minus(j)$
there, so $\mathrm R_N(j)\ge\frac12T_\minus(j)\sum_{m\le\delta/T_\plus(n)}u^\plus_m(n)\ge c''/D_N$ along
the sequence, for $j$ with $j/D_{n+j}$ near $y$ and $N=n+j$, where $D_{n+j}\sim D_n$ because
$j=o(n)$. Together with (b) this gives $D_NP_N(N-j)\ge\rho(y)+c''-o(1)$, and the local law
fails again.

\emph{Part (a).} By (b) and a Riemann sum, $\liminf_N\Prob((N-A_N)/D_N\in[y_1,y_2])\ge\int_{y_1}^{y_2}\rho$
for all $0<y_1<y_2$. Since $\rho$ is a probability density on $(0,\infty)$, the mass of
$(N-A_N)/D_N$ outside $[y_1,y_2]$ is eventually at most $1-\int_{y_1}^{y_2}\rho+o(1)$, which is
small for $y_1$ small and $y_2$ large. Lower bounds on every compact interval together with
tightness give convergence in law. \qed

\subsection{The ends: proof of part (d)}\label{sec:dom-ends}

\emph{Fixed $j$.} In $\mathrm I_N(j)$ only $m\le j$ occur, and for fixed $m$,
$q^\plus_{m+1}(n)=\Prob(S^\plus_{m+1}\ge n)-\Prob(S^\plus_m\ge n)\sim T_\plus(n)$, by
Lemma~\ref{lem:onejump} or by the subexponentiality of a regularly varying tail. So
$\mathrm I_N(j)\sim T_\plus(N)\sum_{m\le j}u^\minus_m(j)=T_\plus(N)u_\minus(j)$. In $\mathrm R_N(j)$ only
$m\le j+1$ occur, and for fixed $m$, $u^\plus_m(n)=o(T_\plus(n))$, because
$f_\plus(x)=T_\plus(x)-T_\plus(x+1)=o(T_\plus(x))$ under (RV), and for $m\ge2$ one of the summands
is at least $n/m$. This uses (RV) for the plus sign only. The minus end is the mirror, with
$q^\minus_m(N-i)\sim T_\minus(N)$ and $u^\minus_m(N-i)=o(T_\minus(N))$ for fixed $m$, and the start
corrections are those of~\eqref{eq:kappa-fresh}.

\emph{The range $1\ll j\ll D_N$.} We show that $P_N(N-j)=T_\plus(N)u_\minus(j)(1+o(1))$ there, under
the condition of (c), $\alpha_\minus\ne1/2$, and $u_\minus\le C\omega$, which holds by (SRT$_\minus$)
when $\alpha_\minus<1$ and trivially when $\mu_\minus<\infty$. For the main term, every $m$ that
carries the mass of $u_\minus(j)$ has $mT_\plus(N)\to0$: if $\mu_\minus<\infty$ the mass sits at
$m\approx j/\mu_\minus\ll1/T_\plus(N)$ by Lemma~\ref{lem:R}(i), and if $\alpha_\minus<1$ the mass
beyond $m=K/T_\minus(j)$ is at most $CK^{1-1/\alpha_\minus+\epsilon}u_\minus(j)$ by Lemma~\ref{lem:R}(ii) and
the lower half of (SRT$_\minus$), while $m\le K/T_\minus(j)$ gives $mT_\plus(N)\le KT_\plus(N)/T_\minus(j)\to0$
because $T_\minus(j)\gg T_\minus(D_N)\sim T_\plus(N)$. Lemma~\ref{lem:Q}(ii) at $t\to0$ gives
$q^\plus_{m+1}(n)=T_\plus(N)(1+o(1))$ on that range, and Lemma~\ref{lem:key} bounds the rest, so
$\mathrm I_N(j)=T_\plus(N)u_\minus(j)(1+O(K^{1-1/\alpha_\minus+\epsilon})+o(1))$, and $K$ is arbitrary.

For the second term with $\alpha_\minus<1$, Lemma~\ref{lem:decay} gives
$q^\minus_m(j)\le CT_\minus(j)e^{-cmT_\minus(j)}$. For $m>\delta/T_\plus(n)$ we have $u^\plus_m(n)\le C_\delta/N$,
so this part is at most $C_\delta/N$, and relative to $T_\plus(N)\omega(j)=T_\plus(N)/(jT_\minus(j))$ it
is at most $CjT_\minus(j)/(NT_\plus(N))\le CD_N/N\to0$, since $xT_\minus(x)$ has a positive index. For
$m\le\delta/T_\plus(n)$ use $u^\plus_m(n)\le C(m/a_\plus(m))T_\plus(N)$. If $\alpha_\plus\ge1/2$, then
$\sum_m(m/a_\plus(m))e^{-cmT_\minus(j)}\le C\lambda(M)M^2/a_\plus(M)$ with $M=1/T_\minus(j)$ and $\lambda$
slowly varying, so the ratio to $T_\plus(N)\omega(j)$ is at most
$C\lambda(1/T_\minus(j))\,j/a_\plus(1/T_\minus(j))$, which tends to $0$. Indeed
$T_\plus(a_\plus(1/T_\minus(j)))\sim T_\minus(j)=o(T_\plus(j))$, so $a_\plus(1/T_\minus(j))/j\to\infty$ by (R2). If
$\alpha_\plus>1/2$ then $\lambda\equiv1$. If $\alpha_\plus=1/2$ then $\alpha_\minus>1/2$, since
$\alpha_\minus\ne1/2$, and $a_\plus(1/T_\minus(j))/j$ is regularly varying in $j$ with the positive index
$2\alpha_\minus-1$, which beats the slowly varying $\lambda$. If $\alpha_\plus<1/2$ the series
$\sum_mm/a_\plus(m)$ converges, the part is at most $CT_\minus(j)T_\plus(N)$, and the ratio is at most
$CjT_\minus(j)^2$, which tends to $0$ when $\alpha_\minus>1/2$. When $\alpha_\minus<1/2$ we use (W$_\plus$)
instead: the part is at most $CT_\minus(j)\,o(1/(D_NT_\plus(N)))$, and the ratio is at most
$o(1)\,jT_\minus(j)^2/(D_NT_\minus(D_N)^2)\le o(1)$, because $xT_\minus(x)^2$ has the positive index
$1-2\alpha_\minus$ and $j\le D_N$~\cite[Theorem~1.5.3]{bgt1989}. At $\alpha_\minus=1/2$ both routes
leave a ratio of slowly varying factors that need not be bounded, which is why we exclude the
line, except that under (SRT$_\plus$) with $\alpha_\plus<1/2$ the ratio is at most
$C\ell_\minus(j)^2/(NT_\plus(N)^2)$ with $T_\minus(x)=x^{-1/2}\ell_\minus(x)$, and $NT_\plus(N)^2$ grows
polynomially. For the second term with $\mu_\minus<\infty$, the part $m>\delta j$ is at most
$C_\delta T_\plus(N)j/a_\plus(j)=o(T_\plus(N))$, since $u^\plus_m(n)\le C(m/a_\plus(m))T_\plus(n)$ and
$m/a_\plus(m)$ has a negative index, and the part $m\le\delta j$ is $o(T_\plus(N))$ as $j\to\infty$ by
the count in the proof of (c), which did not use $j\asymp D_N$.

\emph{The matching.} On $1\ll j\ll D_N$ we have $\Psi(j/D_N)\to1$ and $\kappa_\plus(j)=u_\minus(j)$ for
$j\ge1$, so the two previous paragraphs give the statement there. On $j\asymp D_N$ it is (c),
since $T_\plus(N)u_\minus(j)\Psi(j/D_N)=\rho(j/D_N)(1+o(1))/D_N$ in both cases, by (SRT$_\minus$)
and $jT_\minus(j)\sim y^{1-\alpha_\minus}D_NT_\plus(N)$ when $\alpha_\minus<1$, and by
$u_\minus(j)\to1/\mu_\minus$ when $\mu_\minus<\infty$, while $T_\plus(N)\omega(j)\asymp1/D_N$. For bounded
$j$ it is the fixed-$j$ statement with $\Psi(j/D_N)\to1$. A subsequence argument joins the
three ranges. \qed

\begin{remark}[initial runs]\label{rem:dom-initial}
The first-run decomposition applies unchanged. A minus initial run of length $d$ shifts $j$
by $d$ and adds the defect $T_\plus(N)\sum_i\nu_\minus(j-i)u_\minus(i)$ to the main term. When
$\mu_\minus<\infty$ the kernel $u_\minus$ is bounded, so the defect is
$O(T_\plus(N)\Prob(D\in[j/2,j]))=o(1/D_N)$ for every initial run that is finite almost surely, and
the bulk of the mixed region needs no condition on a minus initial run, as in the finite-mean
case~\cite{paperF}. When $\alpha_\minus<1$ the defect is that of Theorem~\ref{thm:initial} with
$D_NT_\plus(N)$ in place of $NT_\plus(N)$, and the conditions of that theorem on $\nu_\minus$ carry
over. A plus initial run changes the horizon, and we have not written out the uniform bound
for plus initial runs of length close to $N$.
\end{remark}

\ifarxivbuild
\begin{remark}[the big-jump zone toward the lighter end]\label{rem:claimK}
Between the bulk at the scale $D_N$ and the minus end, one macroscopic minus run carries the
deviation, and we expect the two-term kernel of~\cite[Theorem~B]{paperF} with the signs
exchanged and no centering. Assume in addition (LR) for the minus sign. Then, uniformly over
$KD_N\le j\le N-1$ in the double limit $N\to\infty$ then $K\to\infty$, when $\alpha_\minus<1$, we
conjecture that
\[
 P_N(N-j)=\bigl[\gamma_\plus(N-j)\,f_\minus(j)+T_\minus(j)\,u_\plus(N-j)\bigr](1+o(1)),\qquad
 \gamma_\plus(n)=\sum_{r<n}u_\plus(r)\sim\frac{\E E_\plus}{T_\plus(n)},
\]
and the same when $\mu_\minus<\infty$ on $K\lambda_ND_N\le j$ with a factor $\lambda_N$ of order
$(\log N)^{1-\alpha_\plus}$, since then the bulk $\phi$ decays faster than any power. At $j=vN$
and under (SRT$_\plus$) this gives the infinite density
$P_N(N-vN)NT_\plus(N)/T_\minus(N)\to\frac{\sin\pi\alpha_\plus}\pi v^{-\alpha_\minus}(1-v)^{\alpha_\plus}[\alpha_\minus/(\alpha_\plus v)+1/(1-v)]$,
and it reduces to $T_\minus(N)u_\plus(i)$ at $j=N-i$. The exact laws of Section~\ref{sec:verify}
are consistent with the kernel, with ratios that decrease toward $1$ slowly, from $1.66$ at
$v=1/2$ and $N=1000$ to $1.31$ at $N=4000$ in the mixed case $(\alpha_\plus,\alpha_\minus)=(0.6,1.5)$.
A proof should follow the big-jump section of~\cite{paperF}, with a local big-jump bound for
sums of $m$ minus runs that is uniform in $m\le K/T_\plus(n)$, and Lemma~\ref{lem:Q} for the plus
side.
\end{remark}
\fi
